%% notes-tex/figure-3-gate/sections/4-review.tex
%% SOURCE: the ten reviewer reports under review/2026-10-04/ (01 expression before SIMB,
%% 02 expression after SIMB, 03 proteome and joint rounds, 04 baselines and ceilings,
%% 05 architecture and operator, 06 training dynamics, 07 data, splits and metrics,
%% 08 cross-modal evidence and new data, 09 prior reviews and literature, 10 compute and
%% plan). Each report names the results file, config line or W&B run behind every number
%% quoted here. Job ids are from the IGB scheduler on 2026-10-04.

\section{What the audit of 2026-10-04 established\secstatus{todo}}
\label{sec:review}

Ten independent read-only reviews of the expression, proteome and morphology work were run
on 2026-10-04, one scope each, and are filed under \file{review/2026-10-04/}. Where two
reviews reached the same point from different files it is stated once. Every number is
measured unless it is marked as a hypothesis.

\pfstep{No trained model on file beats the best baseline on matched strains} On the test
strains of the model's own partitions, read at the best-validation checkpoint, the graph
transformer scores $0.110$ and $0.153$ on expression (split seeds 0 and 2) against $0.165$
and $0.174$ for a nearest-neighbor average keyed on the interaction graphs, and $0.067$ and
$0.082$ on the proteome against $0.117$ and $0.111$. The model minus the best
validation-selected key is negative in six of six comparisons, by $0.02$ to $0.06$. This
rests on two split seeds per modality, because test predictions were saved for six of forty
checkpoints; validation orders the same way on all four split seeds. Morphology has no
matched comparison at all. A second review reads a different set of checkpoints (v17
reference arm, three seeds per split) at a model edge of about $+0.017$ over the matched
sequence baselines, most of it from split seed 0, so the sign against sequence baselines is
unsettled and the sign against the graph key is not.
\src{review/2026-10-04/04_baselines_and_ceilings.md, 07_data_splits_metrics.md}

\pfstep{Headline numbers that do not hold} The incumbent's $0.196 \pm 0.022$ is the mean
of eight different arms of one round, all at seed 0 on split seed 0, scored by a rolling
maximum; same-seed reruns agree within $0.001$, and no configuration has been replicated at
a long budget. Split seed 0 is a favorable draw: its validation score sits $1.72$ standard
deviations above the mean of twelve baseline draws, and every long run used it. The claim
that the model clears the bilinear baseline by $0.093$ compared a rolling maximum against a
baseline scored on a different permutation of the strains. The rank-64 pair term's
advantage per unit of compute compared one run in each of two projects and fails at
matched budgets. The statistic ``145 runs, correlation of log parameters with score
$-0.129$'' has no generating script and is not reproduced by \file{round_leaderboards.csv},
which gives $-0.001$ over 1{,}182 runs confounded with budget. The manuscript's Results
still carry placeholder correlations of $0.543$ and $0.619$ for expression and morphology.
\src{review/2026-10-04/01_expression_pre_simb.md, 02_expression_post_simb.md,
06_training_dynamics.md, 09_prior_reviews_and_literature.md}

\pfstep{What has a measured effect} Budget, within a single run: $0.147$ at epoch 500,
$0.202$ at 2{,}000, $0.230$ at 4{,}000. The partition: between-partition standard deviation
$0.020$ to $0.030$ against $0.010$ to $0.020$ within, larger than any arm effect. The
masked objective: removing it costs $0.012$ over twelve pairs with a confidence interval
that excludes zero, negative on all four partitions; rounds v19 and v20 run without it.
Embedding content against width-matched random vectors: $+0.068$. Interaction graphs as a
neighbor key against their rewired controls: $+0.150$ on expression and $+0.115$ on the
proteome, four of four split seeds. Among sequence representations ProtT5 alone is at or
above every composite, and the four-part stack the model consumes adds $-0.03$ to $+0.003$
over it. No architecture change has met the campaign's own bar of $+0.02$ on three of four
partitions with the test sign agreeing.
\src{review/2026-10-04/01_expression_pre_simb.md, 02_expression_post_simb.md,
04_baselines_and_ceilings.md}

\pfstep{What the code does at one deletion} The v19 model has 6{,}681{,}548 parameters, of
which 5.69 million, 85 percent, are the first linear layer that projects the 3{,}328
dimensional concatenated embedding; the encoder is 0.59 million and the perturbation
operator 0.10 million. The encoder runs once on the wild-type graph, so its output is the
same for every strain. The operator is a softmax over the perturbed genes; with one deleted
gene its weight is identically 1 and every gene receives the same vector, so its query and
key parameters (16{,}380) receive zero gradient. The Methods describe this softmax; Fig.~1f
draws a weight that differs between genes, which this operator cannot produce at one
deletion. The operator's attention dropout acts on that single weight and drops whole heads
of the deletion signal at random. Nothing marks which gene was deleted: the model predicts
$+0.036$ for the deleted gene's own transcript against a measured $-2.42$, and a predictor
that writes only that one entry scores $0.0175$ per-feature Pearson, the size of the edge in
dispute. The null sink, the sigmoid form that makes the weight depend on the querying gene,
was trained once, for 186 to 189 epochs, with its gate bias near $-4$.
\src{review/2026-10-04/05_architecture_and_operator.md, 07_data_splits_metrics.md,
01_expression_pre_simb.md}

\pfstep{What the training curves show} The runs are neither converged nor memorizing.
Train-side Pearson is still rising at epoch 1{,}200 and flattens below $0.78$ after 10{,}000
epochs. The validation-loss minimum near epoch 70 to 90 sits at the end of a plateau where
predictions are still near-constant, so it marks the launch of the head, not the onset of
overfitting, and the loss-minimum checkpoint is close to a mean predictor. Predictions are
over-dispersed for their correlation (spread ratio about twice the Pearson), which is why
the normalized error exceeds 1; a term that pushes the prediction spread toward the target's
would raise that error further. Wall time does not depend on model size: an epoch costs the
same at 0.68 and 1.45 million parameters, a fifth of it is the eval-mode train pass, and
the encoder is recomputed every step although its output never changes. With the encoder
output cached a CPU step falls from 42~s to 2.6~s, and to 0.23~s on 512 reporter genes.
\src{review/2026-10-04/06_training_dynamics.md, 05_architecture_and_operator.md}

\pfstep{Defects found} Three bear on numbers already reported. The prediction-dump routine
calls the model without the observed-label input, so every dumped validation and test
prediction from the masked-objective lineage comes from a different function than the logged
metric. The v16 expression baselines were taken from the wrong store's partitions. The
per-feature Pearson drops near-constant columns per model, so two arms can be scored on
different feature sets. One bore on running jobs: the conditioned arms were launched without
persistent data-loader workers, the setting whose absence killed four v15 runs; the pending
v20 tasks were resubmitted with it (job 2423297) and tasks 0 and 1 of job 2423190 run as
launched.
\src{review/2026-10-04/05_architecture_and_operator.md, 07_data_splits_metrics.md,
03_proteome_and_joint.md}

\pfstep{What new data can and cannot do inside the window} The local SPELL archive holds
200 studies with genetic perturbations covering 1{,}175 genes, of which 478 are not among
the 1{,}484 Kemmeren deletions: 216 are promoter knockdowns of essential genes, 146 come
from Hughes 2000, and 14 from Hu 2007. On the ridge learning curve 245 to 460 more strains
project to about $+0.004$ to $+0.008$ (hypothesis, one split), below the between-partition
spread. The natural isolates need a reference record, per-allele sequence embeddings that no
builder supports, and a loader fix, five to eight days before a first number. The two
amino-acid screens do not agree with each other because the Cooper screen does not
reproduce its own duplicate strains. No relationship between fitness and the expression or
proteome profile has been measured.
\src{review/2026-10-04/08_cross_modal_and_new_data.md}

\pfstep{The conditioned round} The conditioning step is wired as intended: validation and
test are conditioned, standardization is fit on train, the revealed head leaves the loss
and the metrics. Its comparison against the v19 single-head arms is not like for like (head
slot, loader setting, and an encoder input of values against zeros), so the primary
contrast is each conditioned arm against its permuted control, and the bar is the ridge
from genotype plus the measured modality refit per partition ($0.236$ on expression and
$0.214$ on the proteome on one split), not the unconditioned arm.
\src{review/2026-10-04/03_proteome_and_joint.md}

\section{The plan to the deadline\secstatus{todo}}
\label{sec:plan}

Fewer than twenty days remain. The plan is ordered so that the first week produces a
figure that exists whatever the later experiments show, and so that each later line of work
has a date on which it is dropped.

\pfstep{Days 0 to 4, almost no GPU: the benchmark that cannot fail to exist} Four pieces.
First, the linear and nearest-neighbor baselines for every gene representation on exactly
the strains rounds v19 and v20 use, twelve split seeds, both labels (IGB jobs 2423298 to
2423302, running). Second, one scoring harness applied to everything already trained:
predictions dumped through the same function the metric uses, the deleted gene's own entry
blanked, a feature set fixed in advance by replicate reliability, window means in place of
rolling maxima, test strains, a strain bootstrap. Third, a harness that tells the candidate
limits apart: a ridge and a bilinear model on 137 to 1{,}100 training strains over five
splits, and the same on planted labels at the measured noise level. Fourth, two gates that
earlier reviews asked for and no one ran: perturbation-derived embeddings of the deleted
gene (its genetic-interaction, chemogenomic, morphology or proteome profile) through the
existing baselines, and measured fitness as an input to the genotype ridge with a shuffled
control. \textbf{Gate, day 4:} the matched table is drawn with ceilings. If the model is
below the graph key on at least three of four split seeds per modality, Figure~3 reports a
benchmark, and the remaining work is to raise the model to the graph key.

\pfstep{Days 2 to 8: many small models} The encoder output is cached, so only the operator
and the heads train, on the both-label store already on IGB. Arms, each against a paired
reference in the same pack: the sigmoid-gated operator with an explicit term in the
reporter and the perturbation and an indicator of the deleted gene; the deleted gene's graph
neighborhood as input; a projection cut from 5.69 million parameters to a frozen
low-dimensional one; ProtT5 alone against the stack; the masked objective on; the perturbed
CLS as the morphology readout once the 30-line port from the interaction branch is in.
Screen at 300 epochs on three split seeds, scored by the mean over epochs 200 to 300.
Hypothesis (untested): the cached model runs at least ten times faster per epoch on a GPU.
\textbf{Gate, day 8:} the operator line is dropped unless an arm beats its paired reference
on at least nine of twelve split seeds at confirmation and reaches the graph key.

\pfstep{Days 8 to 15: confirmation} At most two winners and the reference at full size on
twelve split seeds with a sealed test set, about 2.5 days on three A40s.

\pfstep{The conditioned round, in parallel} It runs to completion on the two \file{cabbi}
cards. \textbf{Gate, about day 10:} with at least six split seeds, a conditioned arm must
exceed its permuted control by $0.02$ and reach the per-partition ridge; otherwise it is
reported as a model that uses a second modality less well than a linear map.

\pfstep{What is not in the window} New deletion loaders from SPELL, unless the day-one
learning curve projects a gain above $0.01$. Natural isolates, unless the dataset-build
machine is back by day 5. The distribution-matching loss term, the dataset token, and any
pooled single-cell screen.

\pfstep{Days 16 to 19} Freeze, read the sealed test set once, draw the figure, write.

\pfstep{Decisions that are the author's} Whether Figure~3 may be a benchmark in which a
graph neighbor average ties or beats the transformer. Whether Fig.~1f or the Methods
changes, since they name different operators. Whether the manuscript's placeholder
correlations and its statement that the model predicts expression and morphology jointly
are removed now. Whether the pending v20 tasks keep both directions.
